\documentclass[10pt,a4paper]{article}
\usepackage[utf8]{inputenc}
\usepackage[T1]{fontenc}
\usepackage[margin=1.4cm,top=1.1cm,bottom=1.1cm]{geometry}
\usepackage{amsmath,amssymb}
\usepackage{enumitem}
\usepackage{array}
\usepackage{setspace}
\usepackage{booktabs}
\setstretch{1.0}
\setlength{\parindent}{0pt}
\setlength{\parskip}{0.15em}
\pagestyle{empty}
\newcommand{\repligne}{\par\vspace{0.42cm}\noindent\dotfill}
\newcommand{\zonecalcul}[1]{\par\vspace{0.15cm}\noindent{\small\emph{Espace de calcul :}}\par\vspace{#1}}
\newcommand{\aide}[1]{\fbox{\begin{minipage}{0.96\textwidth}\small #1\end{minipage}}}
\begin{document}
\begin{center}
\textbf{\large Fiche d'exercices --- Séance 6 : la classification}\\[0.05cm]
{\small Introduction au ML -- Dr. El Hadji Bassirou TOURÉ -- DMI $\cdot$ FST $\cdot$ UCAD \quad $\bullet$ \quad \textbf{Travail à la main, sans machine}}
\end{center}
\vspace{-0.15cm}\hrule\vspace{0.25cm}
{\small Cette fiche reprend, exercice par exercice, les calculs de la Séance 6 : le score logistique passé dans la sigmoïde, la log-loss, la matrice de confusion et ses métriques, le seuil comme curseur, la courbe ROC et l'AUC construites à la main, puis le déséquilibre de classes et le choix de la métrique. Répondre directement sur la fiche.}
\vspace{0.2cm}

\aide{\textbf{Aides numériques (valables pour toute la fiche).} Table de la sigmoïde $\sigma(z)=1/(1+e^{-z})$ :
\begin{center}\renewcommand{\arraystretch}{1.15}\begin{tabular}{c|ccccccc}
$z$ & $-3$ & $-2$ & $-1$ & $0$ & $1$ & $2$ & $3$\\\midrule
$\sigma(z)$ & $0{,}047$ & $0{,}119$ & $0{,}269$ & $0{,}500$ & $0{,}731$ & $0{,}881$ & $0{,}953$\\
\end{tabular}\end{center}
Logarithmes : $\ln 0{,}9=-0{,}105$ ; $\ln 0{,}8=-0{,}223$ ; $\ln 0{,}7=-0{,}357$ ; $\ln 0{,}6=-0{,}511$ ; $\ln 0{,}4=-0{,}916$ ; $\ln 0{,}3=-1{,}204$ ; $\ln 0{,}2=-1{,}609$ ; $\ln 0{,}1=-2{,}303$.}
\vspace{0.4cm}

\textbf{\large Exercice 1 -- Du score à la probabilité, puis à la décision}
\vspace{0.15cm}

Un modèle logistique a appris, sur des variables standardisées (âge, glycémie, fièvre), l'intercept $b=-1{,}2$ et les poids $\mathbf w = (0{,}6;\ -0{,}4;\ 1{,}0)$. Une patiente a pour valeurs standardisées $\mathbf x = (1{,}0;\ 1{,}5;\ 0{,}2)$.

\vspace{0.15cm}
(a) Calculer le score $z = b + \mathbf w\cdot\mathbf x$ terme à terme, puis la probabilité $\hat p = \sigma(z)$ (table).
\zonecalcul{1.2cm}
$z = $ \makebox[2.5cm]{\dotfill} \quad $\hat p = $ \makebox[2.5cm]{\dotfill}

\vspace{0.15cm}
(b) Donner la décision au seuil $0{,}5$, puis au seuil $0{,}2$. \quad seuil $0{,}5$ : $\hat y=$ \makebox[1.2cm]{\dotfill} \quad seuil $0{,}2$ : $\hat y=$ \makebox[1.2cm]{\dotfill}

\vspace{0.15cm}
(c) Le poids de la glycémie est négatif. Que signifie ce signe, « toutes choses égales par ailleurs » ?
\repligne\repligne

\vspace{0.15cm}
(d) Trois autres patients ont les scores $z=+3$, $z=0$, $z=-2$. Pour lequel le modèle est-il le plus incertain ? Pourquoi ?
\repligne

\vspace{0.45cm}
\textbf{\large Exercice 2 -- La log-loss à la main}
\vspace{0.15cm}

Quatre prédictions sont jugées par la log-loss $J = -\frac{1}{n}\sum_i\big[y_i\ln\hat p_i + (1-y_i)\ln(1-\hat p_i)\big]$.

\vspace{0.15cm}
{\small \renewcommand{\arraystretch}{1.9}\begin{tabular}{|c|c|c|c|c|}\hline
patient & $\hat p_i$ & $y_i$ & terme qui survit & perte\\\hline
A & $0{,}7$ & $1$ & & \\\hline
B & $0{,}4$ & $1$ & & \\\hline
C & $0{,}2$ & $0$ & & \\\hline
D & $0{,}9$ & $0$ & & \\\hline
\end{tabular}}

\vspace{0.15cm}
(a) Remplir le tableau (quel terme survit, et la perte de chaque patient). \quad (b) Calculer la log-loss moyenne.
\zonecalcul{0.8cm}
$J = $ \makebox[2.5cm]{\dotfill}

\vspace{0.15cm}
(c) Quel patient coûte le plus cher ? Qu'est-ce que la log-loss punit, que l'accuracy ne voit pas ?
\repligne\repligne

\vspace{0.45cm}
\textbf{\large Exercice 3 -- La matrice de confusion}
\vspace{0.15cm}

\aide{Huit patients de dépistage à Kolda, triés par probabilité prédite décroissante (P $=$ paludisme, $y=1$) :
\begin{center}\renewcommand{\arraystretch}{1.2}\begin{tabular}{c|cccccccc}
patient & $1$ & $2$ & $3$ & $4$ & $5$ & $6$ & $7$ & $8$\\\midrule
$\hat p_i$ & $0{,}95$ & $0{,}85$ & $0{,}70$ & $0{,}55$ & $0{,}45$ & $0{,}35$ & $0{,}20$ & $0{,}10$\\
réalité $y_i$ & $1$ & $1$ & $1$ & $0$ & $0$ & $1$ & $1$ & $0$\\
\end{tabular}\end{center}
Ce jeu sert aux Exercices 3, 4 et 5.}

\vspace{0.2cm}
(a) Au seuil $0{,}5$, écrire la décision $\hat y_i$ et la case (TP, FP, FN, TN) de chaque patient.

\vspace{0.1cm}
{\small \renewcommand{\arraystretch}{1.8}\begin{tabular}{|l|c|c|c|c|c|c|c|c|}\hline
patient & $1$ & $2$ & $3$ & $4$ & $5$ & $6$ & $7$ & $8$\\\hline
$\hat y_i$ (seuil $0{,}5$) & & & & & & & & \\\hline
case & & & & & & & & \\\hline
\end{tabular}}

\vspace{0.2cm}
(b) Compléter la matrice, puis calculer accuracy, precision, recall et F1.

\vspace{0.1cm}
\begin{minipage}{0.36\textwidth}{\small \renewcommand{\arraystretch}{1.9}\begin{tabular}{|c|c|c|}\hline
 & prédit $1$ & prédit $0$\\\hline
réel $1$ & TP $=$ & FN $=$\\\hline
réel $0$ & FP $=$ & TN $=$\\\hline
\end{tabular}}\end{minipage}\hfill
\begin{minipage}{0.60\textwidth}
accuracy $=$ \makebox[2.2cm]{\dotfill} \quad precision $=$ \makebox[2.2cm]{\dotfill}\\[0.35cm]
recall $=$ \makebox[2.2cm]{\dotfill} \quad F1 $=$ \makebox[2.2cm]{\dotfill}
\end{minipage}
\zonecalcul{1.0cm}

(c) Traduire la precision et le recall en une phrase chacun, dans la langue d'un agent de santé.
\repligne\repligne

\vspace{0.45cm}
\textbf{\large Exercice 4 -- Le seuil, un curseur}
\vspace{0.15cm}

Même jeu de huit patients (Exercice 3). Le modèle ne change pas ; seul le seuil passe de $0{,}5$ à $0{,}3$.

\vspace{0.15cm}
(a) Quels patients changent de décision ? \makebox[5cm]{\dotfill}

\vspace{0.15cm}
(b) Recalculer TP, FP, FN, TN, puis precision, recall et F1 au seuil $0{,}3$.
\zonecalcul{1.4cm}
TP $=$ \makebox[1cm]{\dotfill} FP $=$ \makebox[1cm]{\dotfill} FN $=$ \makebox[1cm]{\dotfill} TN $=$ \makebox[1cm]{\dotfill} \quad precision $=$ \makebox[1.6cm]{\dotfill} recall $=$ \makebox[1.6cm]{\dotfill} F1 $=$ \makebox[1.6cm]{\dotfill}

\vspace{0.15cm}
(c) Pour une campagne dont la consigne est « ne manquer presque aucun cas », quel seuil préférer ? Quelle métrique accepte-t-on de sacrifier ?
\repligne\repligne

\vspace{0.45cm}
\textbf{\large Exercice 5 -- Construire la ROC et calculer l'AUC à la main}
\vspace{0.15cm}

Même jeu (Exercice 3) : $5$ positifs, $3$ négatifs. On abaisse le seuil juste sous chaque $\hat p_i$, dans l'ordre décroissant.

\vspace{0.15cm}
(a) Compléter le tableau : chaque positif fait \emph{monter}, chaque négatif pousse \emph{à droite}.

\vspace{0.1cm}
{\small \renewcommand{\arraystretch}{1.7}\begin{tabular}{|c|c|c|c|c|c|c|c|c|c|}\hline
le seuil passe sous & départ & $0{,}95$ & $0{,}85$ & $0{,}70$ & $0{,}55$ & $0{,}45$ & $0{,}35$ & $0{,}20$ & $0{,}10$\\\hline
TPR (sur $5$) & $0$ & & & & & & & & \\\hline
FPR (sur $3$) & $0$ & & & & & & & & \\\hline
\end{tabular}}

\vspace{0.2cm}
(b) Tracer la courbe (escalier) dans le repère ci-dessous.
\begin{center}\setlength{\unitlength}{0.4cm}\begin{picture}(12,12)
\thinlines
\multiput(4,1)(3,0){3}{\line(0,1){10}}
\multiput(1,3)(0,2){5}{\line(1,0){9}}
\thicklines
\put(1,1){\vector(1,0){10}}\put(1,1){\vector(0,1){10.8}}
\put(0.4,0.2){\small $0$}\put(3.6,0.2){\small $\tfrac13$}\put(6.6,0.2){\small $\tfrac23$}\put(9.8,0.2){\small $1$}
\put(-0.1,2.8){\small $0{,}2$}\put(-0.1,4.8){\small $0{,}4$}\put(-0.1,6.8){\small $0{,}6$}\put(-0.1,8.8){\small $0{,}8$}\put(0.2,10.8){\small $1$}
\put(10.6,1.3){\small FPR}\put(1.3,11.4){\small TPR}
\end{picture}\end{center}

(c) Calculer l'AUC par les paires : nombre de paires (positif, négatif), nombre de paires mal ordonnées (négatif au-dessus d'un positif), puis AUC.
\zonecalcul{1.0cm}
paires $=$ \makebox[1.5cm]{\dotfill} \quad mal ordonnées $=$ \makebox[1.5cm]{\dotfill} \quad AUC $=$ \makebox[2cm]{\dotfill}

\vspace{0.15cm}
(d) L'AUC dépend-elle du seuil ? Que juge-t-elle exactement ?
\repligne

\vspace{0.45cm}
\textbf{\large Exercice 6 -- La classe rare et l'accuracy trompeuse}
\vspace{0.15cm}

À Tambacounda, $1\,000$ patients sont dépistés ; $50$ sont réellement malades.

\vspace{0.15cm}
(a) Un « modèle » répond toujours « sain ». Calculer son accuracy, son recall et sa precision.
\zonecalcul{0.8cm}
accuracy $=$ \makebox[2cm]{\dotfill} \quad recall $=$ \makebox[2cm]{\dotfill} \quad precision $=$ \makebox[2.5cm]{\dotfill}

\vspace{0.15cm}
(b) Un vrai modèle affiche accuracy $0{,}93$, recall $0{,}60$. Est-il pire que le précédent ? Justifier.
\repligne\repligne

\vspace{0.15cm}
(c) Trois modèles affichent AUC $=0{,}50$, $0{,}88$ et $0{,}25$. Interpréter chacun.
\repligne\repligne

\vspace{0.15cm}
(d) Que fait \texttt{class\_weight="balanced"} à l'entraînement, et que ne change-t-il pas au jeu de test ?
\repligne

\vspace{0.45cm}
\textbf{\large Exercice 7 -- Une métrique pour chaque mission}
\vspace{0.15cm}

(a) Pour chaque situation, choisir la métrique à privilégier (precision, recall, F1 ou AUC) et justifier.

\vspace{0.1cm}
{\small \renewcommand{\arraystretch}{1.9}\begin{tabular}{|p{9.5cm}|c|p{3.6cm}|}\hline
situation & métrique & raison\\\hline
Dépistage du paludisme, test de confirmation bon marché. & & \\\hline
Envoi d'une équipe mobile coûteuse à chaque alerte. & & \\\hline
Comparer deux modèles avant d'avoir fixé le seuil. & & \\\hline
\end{tabular}}

\vspace{0.2cm}
(b) Un collègue règle le seuil en choisissant celui qui maximise la F1 \emph{sur le jeu de test}, puis annonce cette F1. Qu'est-ce qui ne va pas ? Comment faire proprement ?
\repligne\repligne\repligne

\end{document}
